\chapterimage{intstudents.pdf}
\chapterimagecaption{\textit{Leiv Eirikson Discovering America
}(1893) by Christian Krohg}

\chapter{Life as an International Student}
Moving to a different country for school is a big step -- and an exciting one. Beyond diving deeper into your field of study, you'll be discovering a new city, a new culture (or two, in Montr\'eal's case), and (for many) a new climate.

This section is all about figuring out how to navigate both the practical and cultural sides of this transition. There's information about relevant immigration documents, resources to learn French, the winter season, university-provided health insurance, cultural communities, and scholarships.

\section{New beginnings}
The first step towards starting your graduate program here at McGill is accepting your offer of admission on the application portal. You will also have received \textbf{two official offer letters}-- one from the McGill Graduate and Post-doctoral Studies department and one from the Physics department. The latter will detail the program type, duration, stipend amount, etc.

Next, you will have to upload various documents on the portal verifying your status and confirming your eligibility to study in Qu\'ebec. These include -- an official final transcript from your undergraduate university and immigration documents such as a valid CAQ and study permit (see below). Usually, getting these takes a few weeks and you can submit them in the months leading up to the start date of your program (i.e. June to August for a Fall term start, and October to December for a Winter term start).

\subsection{Papers, Please}
To study in Canada, and specifically Qu\'ebec, you will need the following documents,
\begin{plaindefinition}
\begin{itemize}
    \item a valid \textbf{Passport}
    \item a \textbf{Certificat d'acceptation du Québec (CAQ)} -- issued by MIFI
    \item a \textbf{study permit} -- issued by IRCC
    \item a \textbf{Temporary Resident Visa (TRV)} or \textbf{Electronic Travel Authorization (eTA)} -- issued by IRCC, depends on your citizenship
\end{itemize}
\end{plaindefinition}
The CAQ is a provincial document and the study permit is a federal document. 

You would need both to study at McGill. You will first have to apply for a CAQ, via the \href{https://arrima.immigration-quebec.gouv.qc.ca/aiguillage}{Arrima portal}. You will have to upload a copy of your valid passport, your offer letter, proof of financial support (if you have not received it, you can ask the physics graduate program coordinator for this), and sign a \textit{Declarations, Undertakings, and Authorizations} form. The CAQ application fee is $\$135$ (as of 2026).

Once your CAQ application is approved, you will receive a Provincial Attestation Letter (PAL) confirming the issuance of your CAQ. This means you have a confirmed place in the Qu\'ebec share of study permit applications for the year.

You must then apply for a study permit via the \href{https://www.canada.ca/en/immigration-refugees-citizenship.html}{IRCC portal}. For this, you will mostly need the same documents, with the addition of a PAL, and depending on your specific situation, documents required from a personal checklist created by the application portal. The study permit application fee is $\$150$ (as of 2026). When applying for the study permit, you will also have the option to apply for a Temporary Resident Visa (or renew it, if you already have it and it's about to expire). The fee for this is $\$246.25$ (as of 2026).

You can find more information related to immigration and work/study permits on the \href{https://www.mcgill.ca/internationalstudents/work}{International Student Services (ISS)} website.

\subsection*{Maintained Status}\index{Maintained Status}\index{Deregistration}
If your study permit has expired or is expiring soon, you should apply for maintained status if you wish to continue working as a teaching assistant in the department. More details are available \href{https://www.mcgill.ca/apo/staff-guides/immigration/maintained-status}{here.}

To summarize the website, you need to send a copy of your expiring study permit and evidence that you applied to renew your study permit before the expiry of the old one to \href{mailto:immigration.apo@mcgill.ca}{immigration.apo@mcgill.ca}, while CCing Sam. Once your maintained status is approved, forward the email to Kim, who's in charge of hiring.

In case you are dealing with the risk of deregistration due to expired or invalidated documents, get in contact with the Graduate Studies Advisor Nik Provatas and the Graduate Studies Coordinator Sam Maroussis immediately to navigate the situation.

\subsection*{Fast-Tracking}\index{Fast-Tracking}
Note that the CAQ is a title-dependent document, this means that it is only valid if the position of the student as stated on the transcript (i.e. MSc or PhD) matches the position granted by the CAQ. Therefore, even if a CAQ is set to expire in a given date, if a student changes title before that date is reached it will cause the CAQ to immediately become invalid. This means that in case an international student fast-tracks, they must make sure to apply for a new CAQ immediately so as to not be in violation of Québec policy. On the other hand, study permits are not title-dependent and remain unaffected by changes in program or title. In sum, if an international student is planning to fast-track, they must make sure to apply for a new CAQ before the beginning of their first semester as a PhD student, and apply for a new study permit once the new CAQ is approved to account for the additional time of study.

\subsection{Where My People At?}
Finding your community is one of the most important parts of settling in and Montr\'eal makes it easier than you might expect. McGill hosts dozens of cultural and regional student associations, from country- and region-specific clubs to broader international student groups, many of which host welcome events, dinners, and celebrations tied to major holidays from around the world. These groups are often the fastest way to meet people who understand exactly what you're going through, whether that's missing food from home or navigating a new academic system for the first time.

You can find the full list through the Students' Society of McGill University (SSMU) club directory, and most groups are actively looking for new members at the start of the year. 

The ISS holds \href{https://www.mcgill.ca/internationalstudents/pre-arrival/orientation}{international student orientation}, which includes events like chill hangouts in the student lounge (Brown building), welcome reception, and even a soccer tournament. They also host mini-excursions, taking participants to explore iconic Montréal districts and sample local food and culture.

You can also sign up for ISS's \href{https://www.mcgill.ca/internationalstudents/pre-arrival/first-friends}{First Friends} program, which connects you with other incoming students, or the \href{https://www.mcgill.ca/internationalstudents/pre-arrival/buddy}{New Student Buddy} program, which connects new students with returning, experienced McGill students.


\subsection{Mo' Money, Mo' Scholarships}
Please refer to the section on Scholarships (section \hyperref[section:scholarships]{\textcolor{purble}{\ref*{section:scholarships}}}.) for most of the details on this topic.

For international students, the main point to keep in mind are the following: as of the 2026/27 session, you are eligible for the FRQNT scholarship (both masters and doctoral) and the NSERC doctoral scholarship (but not masters).

As a general advice, apply to every scholarship you come across as long as you are eligible. You never know what pays off. More money is always helpful, plus being able to say you won that award holds weight on its own. Being an FRQ/NSERC scholar is a well-regarded achievement.

\subsection{Winters Go Brrr ... } \index{Winters}
Montr\'eal gets cold. The average physical temperature in the winter is around $-5^\circ$ C to $-15^\circ$ C. You should expect frequent $-20^\circ$ C to $-25^\circ$ C or even lower during cold snaps (think up to $-30^\circ$ C). Wind chill brings the average to somewhere between $-10^\circ$ C and $-25^\circ$ C, with cold spells in the $-30^\circ$ C to $-40^\circ$ C range.

This makes \textbf{winter gear} essential. Give yourself time to adjust to the climate and find the right winter coat, waterproof boots, and layers (you should be ready for 10 to 15 degrees of temperature variation during a single day). Take a local friend's advice. Ideally, begin shopping well before the cold sets in. For deals on cheap winter clothings visit Marshalls or Winners.

Keep an eye out for weather warnings (snowstorms and the like), and also for black ice on the streets (this is a very common slipping hazard). Freezing rain usually comes in early or late in winter when the temperatures oscillate around freezing, and is especially dangerous!

Buy a transportation pass (STM monthly pass for example, which comes at a reduced price for students with the student OPUS card). Biking becomes a lot tougher in the winter, but many bike lanes are cleared of snow at the same time as the street is (or soon after). Montr\'eal also has an underground network pathway, called RÉSO or the underground city. This connects many buildings, venues, metro stations, residential complexes, and universities downtown. Use these if commuting above ground becomes challenging due to the cold!
    
\section{Français, SVP} \index{French}\label{section:french}
McGill is an English-language institution, but Montr\'eal is a predominantly French-speaking city. You can absolutely get by day-to-day with English, especially on campus, but a bit of French goes a long way for everyday life -- ordering at a café, reading signage, or navigating city services. \href{https://www.youtube.com/watch?v=ZnCyqeSSWI8}{Here's} a rather short video to get you up to speed to Qu\'eb\'ecois lingo.

There's a bunch of different ways in which you can learn French at McGill.

\begin{itemize}[leftmargin=10pt]
    \item \href{https://www.mcgill.ca/flc/courses-and-programs}{\textbf{French Language Centre (FLC)}} offers courses a variety of courses, categorized into streams based on your end goal. There's French For General Communication courses (code FRSL, held during daytime, for-credit) which are intended for day-to-day communication situations, and French for Professional Communication courses (code FRPC, held in evenings or on Saturdays, for-credit) for those interested in gaining the confidence to work in French in the Qu\'ebec professional context. \vspace{5pt}    
    
    If you want to go even further in your French-learning journey, you can complete an 18-credit minor which combines language, literature, and critical thinking.\vspace{5pt}    
    
    Students in eligible degree programs can apply for an award to cover the cost of tuition and copyright fee for ONE 3 credit course (\$310) offered by the McGill French Language Centre. More details can be found \href{https://www.mcgill.ca/skillsets/files/skillsets/gps_award_for_french_classes.pdf}{here}.\\

    \item \href{https://pgss.mcgill.ca/en/courses}{\textbf{Post-Graduate Students Society (PGSS)}} also offers French courses in the Fall, Winter, and Summer. These are described as leisure courses (taken not-for-credit) with no examinations. The classes are more conversational and move at a steady pace. They are offered in a hybrid format (in-person at Thomson House with simultaneous virtual access via Zoom). The courses are usually charged at $\$85$. Pro-tip: these fill out very quickly, so be sure to register as soon as the window opens. \\

    \item The \href{https://www.quebec.ca/en/education/learn-french}{\textbf{Qu\'ebec government}} also offers free full-time (25-30 hours of classes per week) and part-time (options for 6/9/12/16/20 hours of classes per week) French courses for immigrants. These usually a bit more intense and are intended to quickly bring you up to an intermediate level of proficiency in French. \\

    \item  \textbf{MGAPS} has been organizing French lunches on Fridays, where people meet in the staff lounge and practice conversational skills. These are meant to be very informal and the attendance is usually a mix of beginner and advanced speakers. Plus, being out and about in the city is usually very helpful to pick up new words, learn slang, and improve speaking skills in a real-world setting.\\

    \item Another excellent option is the \textbf{Summer French Immersion} program hosted by the \href{https://www.mcgill.ca/internationalstudents/once-here/summer-french-immersion-program}{\textbf{ISS}}, where students undertake experiential learning at the Université Laval in Québec City. This includes 18 hours of classes each week and cultural activities (cooking, improv, film, etc.). Accepted students also receive an award covering the cost of tuition and accommodation in student residence (but not food or transportation).     
\end{itemize}

In any case, if the academic year gets too busy, you can always grind Duolingo (not a sponsored message).
 









\section{In Sickness and In Health} \index{Health Insurance!International Health Insurance}\label{section:inthealth}
International students receive health insurance (IHI) through the university, via the insurance provider \textbf{Medavie Blue Cross} (as of 2026). The McGill \href{https://www.mcgill.ca/internationalstudents/health}{International Students Services (ISS) website} has a section pertaining to details about insurance coverage and benefits. You can also call them for any inquiries (see below for contact).

There is also a \href{https://www.mcgill.ca/internationalstudents/health/benefits/handbook}{handbook} containing exact coverage amounts for each type of health service. You can also find leisure travel health benefits here. Students pay for IHI via an auxiliary fee (see e-Billings menu on Minerva, in Account Summary by Term) each semester.

Since the IHI does not cover dental services, PGSS provides a dental insurance (also charged as an auxiliary fee). All international students enrolled in the PGSS Dental Plan also have access to the Gender Affirming Care Plan. Learn more about the Gender Affirming Care Plan \href{https://studentcare.ca/rte/en/McGillUniversitygraduatestudentsPGSS_LGBTQIA2SSupport}{here} and \href{https://studentcare.ca/RTEContent/Document/GAC/GAC-Eng_1SEPT2025.pdf}{here}. 

If you'd like, you can opt out of the PGSS Dental Plan during the change of  \href{https://studentcare.ca/rte/en/McGillUniversitygraduatestudentsPGSS_ChangeofCoverage_ChangeofCoveragePeriod}{coverage period}. To opt out, follow the guide from Studentcare \href{https://studentcare.ca/rte/en/ProofofCoverage}{here}.

\begin{plaincorollary}
    The Medavie Blue Cross app is a great platform to submit and track insurance claims, and also has simple look-up tool for coverages and benefits sorted by type of service.
\end{plaincorollary}

You must activate your health insurance on Minerva (takes < 2 mins.) by going to the Student menu, to the International Health Insurance (IHI) Form - Student. Here, you can confirm your coverage, add dependents, and also request exemption from specific services. You can also print your IHI card (you are usually asked to show this at appointments, on phone is fine) from this section.

The \href{https://www.mcgill.ca/wellness-hub/}{McGill Wellness Hub} is your first point of contact for most health services and questions. You can call them to book appointments with general practitioners, counsellors, nurse, or dietitian directly by calling them (see below for contact). Their telephone lines are open Monday to Friday, from 9:00am to 12:00pm, and from 1:30pm to 2:30pm. 




\section{Useful Contacts}
\begin{itemize}
    \item International Student Services (ISS) \\
    Brown Student Services Building, 3600 McTavish Street, Suite 4100 \\
    Phone: +1 514-398-4349
    
    \item McGill Wellness Hub \\
    Brown Student Services Building, 1070 Avenue Dr. Penfield, Suite 3400 \\
    Phone: +1 514-398-6017


\end{itemize}